Cieľom projektu je porovnať trénovanie vybraných architektúr konvolučných
neurónových sietí, vizuálnych transformerov a hybridných modelov pri riešení
úlohy klasifikácie obrazu pomocou transferového učenia (transfer learning).
Táto dokumentácia pokrýva prvú fázu projektu -- model \textbf{M1} (CNN
ResNet50 s transferovým učením) a dva povinné doplnkové experimenty, ktoré
zadanie vyžaduje aspoň pre jeden model:

\begin{itemize}
    \item vplyv augmentácie trénovacích dát na úspešnosť klasifikácie,
    \item vplyv postupného rozmrazovania vrstiev (\emph{gradual unfreezing})
          predtrénovanej siete oproti trénovaniu iba novej výstupnej vrstvy.
\end{itemize}

Ako dataset bol použitý výber z databázy \textbf{Food-101} -- klasifikácia
obrazov jedál. Z pôvodných 101 tried bolo vybraných 15 najviac vizuálne
odlíšiteľných tried (napr. dumplings, lasagna, pizza, steak, waffles), každá
s 1000 obrázkami, spolu 15\,000 obrázkov. Dataset bol náhodne rozdelený v
pomere 80/10/10 na trénovaciu, validačnú a testovaciu množinu (12\,000/1\,500/
1\,500 obrázkov), pri zachovaní rovnakého zastúpenia tried vo všetkých troch
množinách.

Pre model M1 boli spustené tri nezávislé trénovania (rovnako ako pri
doplnkových experimentoch) a výsledky boli štatisticky vyhodnotené
(priemer a smerodajná odchýlka). Pre najlepší beh je zobrazený priebeh
trénovania (Loss a Accuracy na trénovacích a validačných dátach), na
testovacích dátach bola vyčíslená presnosť (Accuracy), presnosť (Precision)
a senzitivita (Recall), spolu s kontingenčnou tabuľkou. Pre vybrané testovacie
vzorky bola vykonaná vizualizácia rozhodovania siete metódou GradCAM.
